\documentclass{article}
\usepackage[utf8]{inputenc}
\usepackage{amsmath}
\usepackage{amssymb}
\usepackage{booktabs}
\usepackage{geometry}
\geometry{a4paper, margin=1in}
\usepackage[margin=1in]{geometry}
\usepackage{amsmath, amssymb, amsthm}
\usepackage{booktabs}
\usepackage{array}
\usepackage{xcolor}
\usepackage{listings}
\usepackage{hyperref}
\usepackage{enumitem}
\usepackage{fancyhdr}
\usepackage{titlesec}
\usepackage{multirow}
\usepackage{tabularx}
\usepackage{tcolorbox}
\tcbuselibrary{skins, breakable}

% Colors
\definecolor{codeblue}{RGB}{30, 100, 180}
\definecolor{codegray}{RGB}{245, 245, 245}
\definecolor{sectionblue}{RGB}{31, 56, 100}
\definecolor{bugred}{RGB}{180, 30, 30}
\definecolor{okgreen}{RGB}{0, 120, 60}
\definecolor{warnyellow}{RGB}{200, 140, 0}
\definecolor{tableheader}{RGB}{213, 232, 240}

% Section formatting
\titleformat{\section}{\large\bfseries\color{sectionblue}}{\thesection}{1em}{}
\titleformat{\subsection}{\normalsize\bfseries\color{codeblue}}{\thesubsection}{1em}{}
\titleformat{\subsubsection}{\normalsize\bfseries}{\thesubsubsection}{1em}{}

% Code block style
\lstset{
  backgroundcolor=\color{codegray},
  basicstyle=\ttfamily\small,
  breaklines=true,
  frame=single,
  rulecolor=\color{gray!40},
  framesep=5pt,
  xleftmargin=10pt,
  xrightmargin=10pt,
  showstringspaces=false,
  columns=fullflexible,
  keepspaces=true
}

% Colored boxes
\tcbset{
  bugbox/.style={
    colback=red!5, colframe=bugred, fonttitle=\bfseries,
    breakable, left=6pt, right=6pt, top=4pt, bottom=4pt
  },
  infobox/.style={
    colback=blue!5, colframe=codeblue, fonttitle=\bfseries,
    breakable, left=6pt, right=6pt, top=4pt, bottom=4pt
  },
  greenbox/.style={
    colback=green!5, colframe=okgreen, fonttitle=\bfseries,
    breakable, left=6pt, right=6pt, top=4pt, bottom=4pt
  },
  warnbox/.style={
    colback=yellow!5, colframe=warnyellow, fonttitle=\bfseries,
    breakable, left=6pt, right=6pt, top=4pt, bottom=4pt
  }
}

\title{One-shot Tiling and Fusion LP Formulation\\for Domain-Optimized Deep Learning Accelerators V1}
\author{CORGI Solver Systems}
\date{March 2026}

\begin{document}

\maketitle

\section{Overview between V0 and V1}

V0 solves fusion and rematerialization jointly, with tiling decisions provided externally by Timeloop
as fixed constants $T_i^{\min}$, $T_i^{\max}$, $t_i^k$. Observation: the best tile size for an operation depends on what tensors the fusion plan
can pin, and the best fusion plan depends on the working-set sizes that tiling produces.
Sequential optimization misses trade-off tension of the form:


Accept a tiling configuration $c'$ that reduces compute efficiency by 5\% relative to
Timeloop's greedy choice, but shrinks the activation working set enough to fit an additional
weight tensor in the Global Buffer, removing the DRAM bottleneck entirely and yielding
a net 20\%+ cycle reduction.


This document if for V1: formulates tiling and fusion as a single LP solved by the CORGIsolver, replacing the sequential Timeloop - FAST Fusion ILP pipeline
for the inner optimization. Vizier's outer loop over hardware configurations is unchanged.

\section{Pipeline Mapping}

The outer pipeline structure is preserved. The joint LP replaces steps 2 and 3 of the sequential pipeline.

\begin{table}[h!]
\centering
\begin{tabular}{llll}
\toprule
\textbf{Component} & \textbf{Original FAST} & \textbf{MultiPeriod V0} & \textbf{V1 (This doc)} \\
\midrule
Vizier            & Hardware config & Hardware config & Hardware config \\
Timeloop          & Single run per op & Single run per op & Run per op \textit{per config} \\
Tiling vars       & None (Timeloop decides) & None & $y_{i,c}$ (new) \\
Placement vars    & $p_i^k$ (static) & $p_{e,t}^k$ (time-indexed) & $p_{e,t},\, p_{i,t}^W$ (unchanged) \\
Remat vars        & None & $r_{i,t}$ & $r_{i,t}$ (unchanged) \\
Linearization vars & None & None & $w_{i,c,k}$ (new) \\
Joint opt.        & No & No & \textbf{Yes} \\
Skip connections  & Restricted & Restricted & \textbf{Relaxed} \\
\bottomrule
\end{tabular}
\caption{Structural comparison across formulation generations.}
\end{table}

\subsection{What Timeloop Does in V1}

Timeloop is run \textit{offline} for each (op $i$, config $c$) pair in the Pareto-pruned
config set $\mathcal{C}_i$. For each pair it returns the triple of constants used in the LP:

\[
  \bigl(T_i^{\min}(c),\;\; T_i^{\max}(c),\;\; t_i^k(c)\bigr)
  \qquad k \in \{I, O, W\}
\]

where $T_i^{\max}(c)$ is the cycle count when all tensors stream from DRAM under config $c$,
$T_i^{\min}(c)$ is the cycle count when all tensors are on-chip under config $c$, and
$t_i^k(c)$ is the cycle saving from pinning tensor of type $k$ in the Global Buffer under
config $c$. These are precomputed constants; the LP solver sees no Timeloop calls at solve
time.

\section{Pareto Pruning Pre-Pass}

The raw tiling search space for a 7-dimensional convolution loop with 3 memory levels is
$O(10^{7})$ configurations per operation. Direct enumeration is infeasible for precomputation?
We reduce to a tractable Pareto-optimal subset $\mathcal{C}_i$ for each op $i$ before
the LP is constructed.

\subsection{Pruning Criteria}

A configuration $c$ is retained in $\mathcal{C}_i$ \textbf{if and only if:}

\begin{enumerate}[noitemsep]
  \item \textbf{L1 mem feasibility:} the tiled working set fits in the L1 scratchpad capacity
        $C_{L1}$. Infeasible configs are dropped before Pareto evaluation.
  \item \textbf{Pareto-optimality} on the two-objective frontier of
        $(T_i^{\max}(c),\; \text{working-set size in Global Buffer under }c)$.
        A config is dominated if another config achieves both lower $T_i^{\max}$ and
        a smaller working set, and is therefore eliminated.
\end{enumerate}


In practice (ref Zhongwei), Pareto pruning reduces $O(10^7)$ raw configs to $|\mathcal{C}_i| \approx 10^2$
to $10^3$ per op. This document sizes the formulation at $|\mathcal{C}_i| = 1{,}000$
as an upper bound; the LP scales linearly with $|\mathcal{C}_i|$.

\subsection{Precomputation Cost? Need to check}

Timeloop must run once per (op, config) pair. For EfficientNet-B7 with $L = 400$ ops and
$|\mathcal{C}_i| = 1{,}000$:

\[
  400 \times 1{,}000 = 400{,}000 \text{ Timeloop evaluations}
\]

At an average of 5 seconds per evaluation, this is approximately 23 hours wall-clock time
on a single machine, or under 30 minutes with 512-way parallelism. Results are cached and
reused across all Vizier trials that share the same hardware configuration class.

\section{Linearization of the Bilinear Interaction}

The joint objective requires the product of two continuous variables: the tiling selection
$y_{i,c} \in [0,1]$ and the tensor pinning indicator $p \in [0,1]$ at the time op $i$
executes ($t = o(i)$). This bilinear term is linearized exactly via the
McCormick substitution.

\subsection{Bilinear Term Structure}

The execution time saving for op $i$ from pinning tensor $k$ under config $c$ is:

\[
  t_i^k(c) \cdot y_{i,c} \cdot p_{\text{assoc}(i,k),\, o(i)}
\]

where $\text{assoc}(i,k)$ maps op $i$ and tensor type $k \in \{I, O, W\}$ to the
corresponding placement variable. We introduce the auxiliary variable:

\[
  w_{i,c,k} \;:=\; y_{i,c} \cdot p_{\text{assoc}(i,k),\, o(i)}
\]

\subsection{McCormick Linearization}

Since both $y_{i,c}$ and $p \in [0,1]$, the product is linearized exactly by four
linear constraints:

\[
\begin{aligned}
  w_{i,c,k} &\leq y_{i,c} \\
  w_{i,c,k} &\leq p_{\text{assoc}(i,k),\, o(i)} \\
  w_{i,c,k} &\geq y_{i,c} + p_{\text{assoc}(i,k),\, o(i)} - 1 \\
  w_{i,c,k} &\geq 0
\end{aligned}
\qquad \forall\, i \in \mathcal{N},\; c \in \mathcal{C}_i,\; k \in \{I, O, W\}
\]

This linearization is \textbf{exact} at integer solutions: when $y_{i,c} \in \{0,1\}$ and
$p \in \{0,1\}$, the four constraints force $w_{i,c,k} = y_{i,c} \cdot p$ exactly.
No big-$M$ constants are required, and the LP relaxation bound is tight.

\section{The Joint V1 LP Formulation}

\subsection{Decision Variables}

\[
  y_{i,c} \in [0,1], \quad
  p_{e,t} \in [0,1], \quad
  p_{i,t}^W \in [0,1], \quad
  r_{i,t} \in [0,1], \quad
  w_{i,c,k} \in [0,1], \quad
  T_i \geq 0
\]

In an LP relaxation all variables are continuous on $[0,1]$; integrality is enforced
post-solve via rounding or by re-solving as MIP if the relaxation gap is unacceptable.

\subsection{Full Formulation}

\[
\begin{aligned}
\text{minimize} \quad
  & \sum_{i \in \mathcal{N}} T_i
  \;+\; \sum_{i \in \mathcal{N}} \sum_{t \in \mathcal{T}} c_i^{\text{remat}}\, r_{i,t} \\[6pt]
%
\text{subject to} \\[4pt]
%
% (1) One tiling config per op
  & \sum_{c \in \mathcal{C}_i} y_{i,c} = 1
  & \forall\, i \in \mathcal{N} \\[4pt]
%
% (2) Execution time -- upper bound (joint tiling + fusion)
  & T_i \;\geq\; \sum_{c \in \mathcal{C}_i} T_i^{\max}(c)\, y_{i,c}
    \;-\; \sum_{c \in \mathcal{C}_i} \sum_{k \in \{I,O,W\}} t_i^k(c)\, w_{i,c,k}
  & \forall\, i \in \mathcal{N} \\[4pt]
%
% (3) Execution time -- lower bound (compute-bound floor)
  & T_i \;\geq\; \sum_{c \in \mathcal{C}_i} T_i^{\min}(c)\, y_{i,c}
  & \forall\, i \in \mathcal{N} \\[4pt]
%
% (4) McCormick linearization
  & w_{i,c,k} \leq y_{i,c}
  & \forall\, i,\, c \in \mathcal{C}_i,\, k \\[2pt]
  & w_{i,c,k} \leq p_{\text{assoc}(i,k),\, o(i)}
  & \forall\, i,\, c \in \mathcal{C}_i,\, k \\[2pt]
  & w_{i,c,k} \geq y_{i,c} + p_{\text{assoc}(i,k),\, o(i)} - 1
  & \forall\, i,\, c \in \mathcal{C}_i,\, k \\[4pt]
%
% (5) Global memory capacity -- per timestep
  & \sum_{e \in \mathcal{E}} d_e\, p_{e,t}
    \;+\; \sum_{i \in \mathcal{N}} d_i^W\, p_{i,t}^W
    \;\leq\; C_{\text{GM}}
  & \forall\, t \in \mathcal{T} \\[4pt]
%
% (6) Activation persistence (eviction / remat)
  & p_{e,t} \leq p_{e,t-1} + r_{\text{src}(e),\, t}
  & \forall\, e \in \mathcal{E},\; t > o(\text{src}(e)) \\[4pt]
%
% (7) Weight monotone non-increase (no reload without remat)
  & p_{i,t}^W \leq p_{i,t-1}^W
  & \forall\, i \in \mathcal{N},\; t > o(i) \\[4pt]
%
% (8) Remat only after first execution
  & r_{i,t} = 0
  & \forall\, i \in \mathcal{N},\; t \leq o(i) \\[4pt]
%
% (9) Remat requires input activations to be present
  & r_{i,t} \leq p_{e,t}
  & \forall\, e \in \text{in}(i),\; t > o(i) \\[4pt]
%
% (10) Remat requires weights to be present
  & r_{i,t} \leq p_{i,t}^W
  & \forall\, i \in \mathcal{N},\; t > o(i) \\[4pt]
%
% (11) Boundary: tensors zero before produced
  & p_{e,t} = 0
  & \forall\, e \in \mathcal{E},\; t < o(\text{src}(e)) \\[2pt]
  & p_{i,t}^W = 0
  & \forall\, i \in \mathcal{N},\; t < o(i) \\[4pt]
%
% (12) Non-negativity
  & T_i,\; p_{e,t},\; p_{i,t}^W,\; r_{i,t},\; w_{i,c,k} \geq 0
\end{aligned}
\]

\subsection{Notes on the Constraints}

\begin{itemize}
    \item Constraints (2) and (3) jointly express that the LP must commit to a tiling config
$c$ (via $y_{i,c}$) which sets both the baseline compute cost $T_i^{\max}(c)$ and
the per-tensor DRAM savings $t_i^k(c)$. A config with smaller tiles may have higher
$T_i^{\max}(c)$ (lower compute utilization) but also smaller tensor working sets,
enabling more aggressive pinning. The LP balances these simultaneously rather than
sequentially.
\item Constraint (6) is indexed over edges $e \in \mathcal{E}$ rather than ops.
A node with fan-out degree $d$ generates $d$ edges, each with its own $p_{e,t}$ variable
and persistence constraint. This correctly models the case where one consumer of a
skip connection benefit from on-chip residual while another does not, without the
big-$M$ consecutive-execution restriction of the original FAST fusion ILP.

\end{itemize}


\section{Variable Definitions and Dimension Estimates for V1}

EfficientNet-B7 estimates: $L = 400$ ops, $T = 400$ timesteps,
$|\mathcal{C}_i| = 1{,}000$ Pareto-pruned configs, $|\mathcal{E}| \approx 400$ edges.

\begin{table}[h!]
\centering
\begin{tabular}{lrrl}
\toprule
\textbf{Variable} & \textbf{Dimension} & \textbf{Count} & \textbf{Notes} \\
\midrule
$T_i$             & $L \times 1$              & 400           & Continuous, unchanged \\
$y_{i,c}$         & $L \times |\mathcal{C}|$  & 400,000       & New: tiling config selection \\
$p_{e,t}$         & $|\mathcal{E}| \times T$  & 160,000       & Edge-indexed, time-indexed \\
$p_{i,t}^W$       & $L \times T$              & 160,000       & Weight placement \\
$r_{i,t}$         & $L \times T$              & 160,000       & Rematerialization \\
$w_{i,c,k}$       & $L \times |\mathcal{C}| \times 3$ & 1,200,000 & New: linearization variables \\
\midrule
\textbf{Total}    &                           & \textbf{2,080,400} & \\
\midrule
Constraint matrix $A$ & $\approx 5{,}000{,}000 \times 2{,}080{,}400$ & & \\
Estimated NNZ         & $\approx$ 15M -- 25M   & & $\ll$ CORGI 6.5B limit \\
\bottomrule
\end{tabular}
\caption{Variable count at $L = 400$, $T = 400$, $|\mathcal{C}_i| = 1{,}000$.}
\end{table}

\subsection{NNZ Breakdown by Constraint Block}

\begin{table}[h!]
\centering
\begin{tabular}{lrl}
\toprule
\textbf{Constraint Block} & \textbf{Estimated NNZ} & \textbf{Source} \\
\midrule
One-config-per-op (1)           & 400,000      & $L \times |\mathcal{C}|$ entries \\
Exec.\ time upper bound (2)     & 4,800,000    & $|\mathcal{C}| \times 4$ terms per op \\
Exec.\ time lower bound (3)     & 400,000      & $|\mathcal{C}|$ terms per op \\
McCormick linearization (4)     & 7,200,000    & 3 constraints $\times$ 2 vars $\times$ $L|\mathcal{C}|3$ \\
Global memory capacity (5)      & 320,000      & $T \times (|\mathcal{E}| + L)$ \\
Activation persistence (6)      & 800,000      & $|\mathcal{E}| \times T \times 2$ \\
Weight monotone (7)             & 320,000      & $L \times T \times 2$ \\
Remat feasibility (8--10)       & 640,000      & Several $L \times T$ blocks \\
Boundary conditions (11)        & 800,000      & $|\mathcal{E}| \times T + L \times T$ \\
\midrule
\textbf{Total (approx.)}        & \textbf{$\sim$15.7M} & \\
\bottomrule
\end{tabular}
\caption{NNZ estimate by constraint block. Well within CORGI's capacity.}
\end{table}

\section{Parameter Constants (Inputs Stored from Precomput)}

\begin{table}[h!]
\centering
\begin{tabular}{@{}llp{7.5cm}@{}}
\toprule
\textbf{Constant} & \textbf{Source} & \textbf{Description} \\
\midrule
$T_i^{\max}(c)$   & Timeloop        & Cycle count for op $i$ under config $c$, all tensors in DRAM. \\
$T_i^{\min}(c)$   & Timeloop        & Cycle count for op $i$ under config $c$, all tensors on-chip. \\
$t_i^k(c)$        & Timeloop        & Cycle saving from pinning tensor $k$ in Global Buffer under config $c$. \\
$d_e$             & HLO graph       & Size in bytes of activation tensor on edge $e$. \\
$d_i^W$           & HLO graph       & Size in bytes of weight tensor for op $i$. \\
$c_i^{\text{remat}}$ & Timeloop / analytical & Cycle cost of recomputing op $i$. \\
$C_{\text{GM}}$   & Vizier          & Global Memory capacity in bytes (e.g., 128\,MiB). \\
$o(i)$            & HLO graph       & Pre-determined execution order of fusion regions. \\
$\mathcal{C}_i$   & Pareto pruning  & Set of Pareto-optimal tiling configs for op $i$. \\
\bottomrule
\end{tabular}
\caption{Input constants. All are precomputed before the LP.}
\end{table}

\section{Comparison to V0 and FAST static}

\begin{table}[h!]
\centering
\begin{tabular}{lccc}
\toprule
\textbf{Property} & \textbf{FAST Fusion ILP} & \textbf{Multi-Period V0} & \textbf{This LP (V1)} \\
\midrule
Joint tiling + fusion            & \texttimes & \texttimes & \checkmark \\
Time-indexed placement           & \texttimes & \checkmark & \checkmark \\
Rematerialization                & \texttimes & \checkmark & \checkmark \\
Relaxed skip-connection liveness & \texttimes & Partial    & \checkmark \\
Config-dependent DRAM savings    & \texttimes & \texttimes & \checkmark \\
Bilinear terms                   & None       & None       & Linearized exactly \\
Solver                           & SCIP (CPU) & CORGI      & CORGI \\
Solve time (reported/estimated)  & 20 min-Hours (ref Dan) & Fast     & Est.\ fast \\
\bottomrule
\end{tabular}
\caption{Static, V0, V1 comparison}
\end{table}


\newpage


\[
\begin{aligned}
\text{minimize} \quad
  & \sum_{i \in \mathcal{N}} T_i
  \;+\; \sum_{i \in \mathcal{N}} \sum_{t \in \mathcal{T}} c_i^{\text{remat}}\, r_{i,t} \\[6pt]
%
\text{subject to} \\[4pt]
%
% (1) One tiling config per op
  & \sum_{c \in \mathcal{C}_i} y_{i,c} = 1
  & \forall\, i \in \mathcal{N} \\[4pt]
%
% (2) Execution time -- upper bound (joint tiling + fusion)
  & T_i \;\geq\; \sum_{c \in \mathcal{C}_i} T_i^{\max}(c)\, y_{i,c}
    \;-\; \sum_{c \in \mathcal{C}_i} \sum_{k \in \{I,O,W\}} t_i^k(c)\, w_{i,c,k}
  & \forall\, i \in \mathcal{N} \\[4pt]
%
% (3) Execution time -- lower bound (compute-bound floor)
  & T_i \;\geq\; \sum_{c \in \mathcal{C}_i} T_i^{\min}(c)\, y_{i,c}
  & \forall\, i \in \mathcal{N} \\[4pt]
%
% (4) McCormick linearization
  & w_{i,c,k} \leq y_{i,c}
  & \forall\, i,\, c \in \mathcal{C}_i,\, k \\[2pt]
  & w_{i,c,k} \leq p_{\text{assoc}(i,k),\, o(i)}
  & \forall\, i,\, c \in \mathcal{C}_i,\, k \\[2pt]
  & w_{i,c,k} \geq y_{i,c} + p_{\text{assoc}(i,k),\, o(i)} - 1
  & \forall\, i,\, c \in \mathcal{C}_i,\, k \\[4pt]
%
% (5) Global memory capacity -- per timestep
  & \sum_{e \in \mathcal{E}} d_e\, p_{e,t}
    \;+\; \sum_{i \in \mathcal{N}} d_i^W\, p_{i,t}^W
    \;\leq\; C_{\text{GM}}
  & \forall\, t \in \mathcal{T} \\[4pt]
%
% (6) Activation persistence (eviction / remat)
  & p_{e,t} \leq p_{e,t-1} + r_{\text{src}(e),\, t}
  & \forall\, e \in \mathcal{E},\; t > o(\text{src}(e)) \\[4pt]
%
% (7) Weight monotone non-increase (no reload without remat)
  & p_{i,t}^W \leq p_{i,t-1}^W
  & \forall\, i \in \mathcal{N},\; t > o(i) \\[4pt]
%
% (8) Remat only after first execution
  & r_{i,t} = 0
  & \forall\, i \in \mathcal{N},\; t \leq o(i) \\[4pt]
%
% (9) Remat requires input activations to be present
  & r_{i,t} \leq p_{e,t}
  & \forall\, e \in \text{in}(i),\; t > o(i) \\[4pt]
%
% (10) Remat requires weights to be present
  & r_{i,t} \leq p_{i,t}^W
  & \forall\, i \in \mathcal{N},\; t > o(i) \\[4pt]
%
% (11) Boundary: tensors zero before produced
  & p_{e,t} = 0
  & \forall\, e \in \mathcal{E},\; t < o(\text{src}(e)) \\[2pt]
  & p_{i,t}^W = 0
  & \forall\, i \in \mathcal{N},\; t < o(i) \\[4pt]
%
% (12) Non-negativity
  & T_i,\; p_{e,t},\; p_{i,t}^W,\; r_{i,t},\; w_{i,c,k} \geq 0
\end{aligned}
\]

\newpage
\begin{equation*}
\begin{aligned}
\text{minimize} & \sum_{i\in \mathcal N} T_i+\sum_{i\in\mathcal N}\sum_{t\in\mathcal T} c_{i}^{\text{remat}} r_{i,t}\\
\text{subject to } & T_i\geq T_i^{\text{max}}-\sum_{e\in\text{in}(i)}t_e^I\cdot p_{e,o(i)}-t_i^O\cdot p_{e',o(i)}-t_i^W\cdot p_{i,o(i)}^W\\
& T_i\geq T_i^{\text{min}}\\
& \sum_e d_e p_{e,t}+\sum_i d_i^W\cdot p_{i,t}^W \leq C_{\text{GM}} \quad \forall t\\
& p_{e,t}\leq p_{e,t-1}+r_{\text{src}(e),t}\quad \forall e,t>o(\text{src}(e))\\
& r_{i,t}=0\quad \forall t\leq o(i)\\
& r_{i,t} \leq p_{e,t} \quad \forall e\in\text{in}(i)\quad \forall t>o(i)\\
&r_{i,t}\leq p_{i,t}^W\quad \forall i,\forall t>o(i)\\
& p_{i,t}^W\leq p_{i,t-1}^W\\
% & p_{i,t}^O\leq p_{j,t}^I\\
% & 0\leq p_{i,t}^k,r_{i,t}\leq 1 \\
& 0\leq  p_{e,t}, p^W_{i,t}, r_{i,t}\leq 1\\
& T_i\geq 0 \\
& p_{e,t} = 0 \quad \forall t<o(\text{src}(e)) \\
& p_{i,t}^W=0 \quad \forall t<o(i)
\end{aligned}
\end{equation*}



\end{document}

\end{document}